%! iTeXMac(project): thesis.pTexMac
%! iTeXMac(input):thesis.tex 
\chapter{Introduction}
\label{chap:intro}

One of the most impressive characteristics found in mammals is their ability to adapt to new environments.  This adaptability relies on learning and memory.  The ability to generalise learnt responses to novel situations is of significant evolutionary advantage, allowing the animal to respond to changing environments and to extract the important similarities between situations.   Without the ability to generalise, every situation appears to be unique and atomic, requiring a unique behavioural response.  How is it that animals are able to learn that pressing a red lever will release food regardless of the context, and then later learn that the lever will only work after a bell is rung?

This ability to generalise about the results of actions and to select an appropriate action given a particular situation is vital to the survival of the individual.  The ability to store the results of previous experience is called memory.  The taxonomy of memory we will use is in line with the work of~\citet{Tulving1987} in which he breaks memory into three categories -- \term{procedural}, \term{semantic}, and \term{episodic}.

Procedural memory is the most primitive of the memory types.  This type of memory describes an animal's ability to learn to perform a task, a ``procedure'', allowing the learning of sequences of actions.  Procedural memory does not require an explicit representation of time, it only needs to store the next action to perform.   For example, the ability to hit a tennis ball is a procedural memory.  The ability to hit the ball does not depend on memory of a particular instance in which you hit a tennis ball.  Instead, the memory is stored as a sequence of actions coupled with a perceptual feedback system that keeps the racquet on target. 
   
Semantic memory refers to the learning of facts and concepts, the meaning behind the structural relationships in the environment.  This is the most abstract form of memory.  Such memories do not elicit action, nor are they related to any distinct event.  However, semantic memory is very important when thinking and talking about the world around us.  These memories allow us to extract meaning from repeated experience and to develop theories about the relationships between experiences.  By testing these theories of association we can develop the fundamental structure needed for abstract thought.
 
Episodic memory is memory for specific events or episodes. This requires some form of temporal component to locate the event in time.  This is the type of memory that is most commonly associated with the word ``memory'' and it provides us with our sense of identity, allowing us to feel the passage of time and to reflect on events in the past.  This reflection on the past allows for a certain amount of learning to occur even after an event has passed.

Our memories are precious to each of us, they define who we are and how we respond to our environment. The study of memory and the mechanisms behind memory is critical both as an introspective analysis of ourselves and as a computational problem from a computer science perspective.  By modelling biological memory we hope to learn more about how memory works in mammals and gain insight into how to develop software which can interact with humans in a more natural way.

This thesis investigates just a small part of the picture of memory.  It is focussed on how memory is stored over time and in particular how mammals avoid the problems associated with information overload and the subsequent forgetting of old information.  It is obvious that critical information is encoded early in an animal's life, and that certain responses learnt in these early years need to remain robust throughout the animal's lifetime.  This poses a dilemma regarding how to protect the relevant information from the past while remaining flexible enough to adapt to new environments.

\section{Memory Mechanisms}

The exact mechanism used to store long term memory in biological neural networks is still unknown. We know that there are certain areas of the brain that are vital for the development and storage of information.  It has been shown that both the neocortex and the hippocampus are important for episodic memory~\citep{Atkinson1968,Meeter2003}.  Models of short term memory (STM) and long term memory (LTM) rely heavily on lesion studies, CAT scans, and more recently fMRI, to give us an insight into the functioning of the mammalian brain. These methods investigate whole brain regions, each containing millions of individual neurons,the functional units of the brain, and hundreds of millions of synapses, the connections between the neurons.  The results from these studies show us the areas of the brain that are important, and the structures that are involved in memory formation and recall.  However, these methods are currently still too coarse to detect the minute and subtle changes that are occurring during memory formation.

In 1949, Donald Hebb suggested a mechanism that could alter the synaptic connections between neurons, so that structure would develop in the network of neurons.  Hebb proposed that when two neurons were active at the same time, the likelihood of the same neurons firing at the same time in the future was increased.
\begin{quotation}
``When an axon of cell A is near enough to excite a cell B and repeatedly and persistently takes part in firing it, some growth process or metabolic change takes place in one or both cells such that A's efficiency, as one of the cells firing B, is increased.''~\citet{Hebb49}.
\end{quotation}

The biological basis for this process was described by \citet{Bliss1973}.  This work expanded on the experimental results from work done in 1966 by L{\o}mo~\citep{Lomo2003,Bliss2003}. The long term change in the effectiveness of a neuron to activate another neuron is called \term{Long Term Potentiation} (LTP)~\citep{Bliss1973}. LTP has been further divided into early-phase LTP (E-LTP) and late-phase LTP (L-LTP).  E-LTP seems to last less than an hour~\citep{Mayford1995} and so is unlikely to be responsible for the lifelong changes required for long term memory.  L-LTP however requires new protein synthesis and synaptic growth, and appears to last for years~\citep{Abel1997}.  This research provides a link between the microscopic changes occurring at a synaptic level and the behavioural phenomenon of memory.  It is still not entirely clear whether L-LTP is solely responsible for lifelong memory, but it is likely that it has a significant role to play~\citep{Abraham2000}.  

The hippocampus is a focal point for the study of memory as it seems to be involved in most aspects of information storage, maintenance and retrieval.  Damage to the hippocampus severely diminishes short term memory, and the ability to transform short term memories into long term memories.  The hippocampus has been shown to be highly active when subjects are learning new material~\citep{Davachi2003}. As well as serving an important role in memory the hippocampus has also been identified as important for navigation and learning the relationship between events and locations. This is most obvious in the studies of rat navigation in which rewards affect the firing patterns of cells in the hippocampus~\citep{Burgess94}.

The main feature that distinguishes human memory from computer memory is our ability to discover associations between events and to access information by analogy.  This is extremely difficult with computers as they store information as discrete and isolated patterns of binary digits.  Mammalian memories seem to be overlapping and they share resources with other memories forming interconnected representations.  This overlapping of representation allows mammals to quickly access relationships between similar items and to respond to them in similar ways.  When faced with a deadly carnivore it is a good idea to run away, irrespective of the exact number of teeth, the length of the legs or the position of the sun.  To find similarities between events using a computer, it is necessary to explicitly work though all the combinations of features to detect the overlaps.  This is an extremely complex task both algorithmically and computationally, requiring metrics to decide what is meant by similar and which of the millions of features of an event are relevant.

\section{Content Addressable Memory}

To enable computers to efficiently find associations between stored information and similar current events, the information needs to be stored in a way that allows comparison and association.  Rather than being accessed purely by the location of the memory, it needs to be recalled based on the content.  The popularity of Google\super{tm} is partly due to the natural way of accessing information.  By providing text that is either on a web page or linked to the page, the Google\super{tm} search returns pages that have content related to the search string.  Google Desktop\super{tm} performs an analysis of the computer's hard drive and stores content information in an index that can be searched quickly for relevant words and terms.  These systems allow users to use textual context to access information, but they are extremely limited and require significant computational effort and additional memory.  What is required is a memory that can be directly accessed by contents.

Content addressable memory (CAM) is the ability to access information from just a partial sample of the original content. For example, when trying to remember an event it is common to retrace your steps and present as much of the original context as possible.  Mammals appear to have the same ability to remember behaviour based on contextual information.  When developing a functional model of mammalian memory it is important to replicate this feature. Computer models of CAM have been developing over many years~\citep{Hanlon1966, Hopfield82, Kohonen1987, Pagiamtzis2006}. The model of CAM that is the focus of this thesis is the Hopfield model. This model has a number of desirable features:
\begin{itemize}
	\item A biologically plausible connection update rule, the Hebbian learning rule.
	\item A simple activation function than can be analysed and provide repeatable results.
	\item A large body of research for comparison of results.
\end{itemize}
The majority of this thesis will deal with variants of the Hopfield network (see Section~\ref{sec:hopfieldnetworks}) and analyse the performance of different learning methods applied to this model of CAM.

The term CAM is also applied to hardware which allows fast recovery of information based on partial content~\citep{Schultz1996}.  These systems provide order of magnitude faster access to information than conventional digital memory.  They are used in specialist areas such as telecommunications where speed is far more important that cost or energy efficiency.  Hardware implementations are beyond the scope of this thesis. We will instead be focusing on software which simulates the properties of biological memory.

\remove{In standard computer memory design significant effort is expended to make sure that the individual components stored do not interfere.  The distance between components on silicon wafers are measured in microns.  This is approaching the limit of what can be achieved with a metal substrate as the distances approach the molecular level.  These components must however be kept separate, as it is fundamental to the von Neuman architecture that program instructions do not accidentally interfere with each other while they are stored in memory.  The only component that can make any change in classical computer design is the CPU.  Without the certainty that all changes are made by the CPU, program execution would be a complicated and unstable process.  Memory that changes content without input from the CPU is flagged as corrupt and is used as a signal to replace the hardware.}

\section{Flexible Representation}

The brain is undergoing constant change.  These changes are not just the transient electro-chemical signals that pulse through the brain, but changes to the shape and connectivity of neurons.  There are also changes to the efficacy and speed of connections already present.  Every neuron, and possibly every synapse, is making changes to its structure and behaviour.  When the brain is damaged, it is often able to recover much of its original functionality. Computer memory on the other hand is fixed and inflexible.  When errors occur the information in that area is usually completely lost, indeed even a single fault can cause the entire system to fail.

There are two reasons that we might want to reconcile these fundamentally different systems.  The first is to help us more fully understand the nature of human memory, and the second is to improve the performance of the memory systems in computers, not to make them more accurate, but to make them more useful.

The task that we are particularly interested in is the mammalian ability to change behaviour in response to events that occur over an extended time frame.  Mammalian brains are flexible enough to learn new responses to old situations and retain different behaviours for similar but distinct events.  For example, when a rat is placed in an environment it can learn the location of food sources and is able to navigate to those locations weeks later, even after being exposed to many new environments.

The main problem posed by this type of learning is how to incorporate new information without destroying the memories that have already been stored.  The animal could try to store everything that happens to it, forget those things that were not connected to survival, or just forget all the events and just remember simple relationships.  Given that remembering everything is cognitively expensive, and that mammals seem to remember more than simple relationships, there must be a mechanism that allows smooth integration of new information without loosing too much old information.  

Computer models of neural networks, artificial neural networks (ANNs) suffer from a peculiar problem when presented with new information that needs to be incorporated into an already functioning memory system.  The old information disappears almost immediately new information is learnt \citep{McCloskey1989}.  We will call this catastrophic loss of memory ``catastrophic forgetting'' (CF).  Animals do not seem to suffer from this problem.  What mechanisms are missing from our models of memory that prevent this loss in the face of new learning?  Is it possible to create computer models that are as robust and flexible as the mammalian brain?

We suggest that mammals have solved the problem of catastrophic forgetting by consolidating newly learnt information in an off-line process during sleep.  The role of sleep as part of memory consolidation has been widely explored in the literature with promising results~\citep{Pearlman71, Pearlman73, Pearlman74, Pearlman79, Fishbein81, Smith91, Smith1993, Smith95, Hennevin95, Hennevin1995a, Smith96, Fishbein97, Meeter2004, Muzur05, Stickgold2005}.  
Even with many decades of research the process that protects memory is not fully understood.  
There is even debate about the necessity of sleep for memory protection~\citep{Siegel2001,Vertes2004,Siegel2005} and many questions have arisen regarding the nature of these memory mechanisms. For example, what is the nature of the processing of information, and what might be happening during sleep that allows for improvements in certain types of memory~\citep{Stickgold2001}?

The approach that we will take to investigate some of these questions, is to develop a memory system that shares some of the features of mammalian memory and then to experiment with several different possible mechanisms for preventing the catastrophic forgetting that is common in computational versions of these systems.  The Hopfield network model that we will investigate has two of the key features that we are interested in.  It is a CAM and it suffers from catastrophic forgetting.

\section{Thesis Structure}

In the rest of this thesis we discuss the problem of catastrophic forgetting and look at various solutions to this problem.  Chapter~\ref{chap:Background} gives a background to the problem and the network that will be investigated. Chapter~\ref{chap:CF} describes the catastrophic forgetting problem in multilayer perceptron networks and some of the proposed solutions. Chapter~\ref{chap:SolutionsCF} moves the investigation of catastrophic forgetting into Hopfield networks and describes previous solutions to this class of problems as well as our pseudorehearsal procedure~\citep{Robins1998,Robins1999}. One of the features of pseudorehearsal is the selection of the patterns/memories to rehearse.  Chapter~\ref{chap:distinguish} investigates ways to distinguish between real memories and spurious memories, and presents a new technique which involves the comparisons of the inputs to individual units~\citep{Robins2004}.  This technique is then applied to both the pseudorehearsal process and the unlearning process~\citep{Hopfield83,vanHemmen1997} in Chapter~\ref{chap:Enhancing}, to establish if performance can be improved.  Finally Chapter~\ref{chap:Discussion} discusses future directions for this research, and the link to sleep and dreaming. 

